% ======================================================================
\section{Relecture : sections 3 à 6 de l'article}

Cette section répond, dans l'ordre de l'article, aux questions posées à la relecture des sections 3 (estimateurs d'environnement) à 6 (choix d'analyse).

\subsection{Section 3 : les environnements}

\paragraph{La corrélation de Spearman et la valeur 0.72.} On classe tous les points une fois selon $r$ (ASTRA) et une fois selon $\delta_R$ (density split), puis on calcule la corrélation entre les deux classements. Elle vaut 1 si les deux méthodes ordonnent les points exactement de la même façon, 0 si leurs ordres n'ont rien à voir. On utilise Spearman (sur les rangs) plutôt que la corrélation habituelle parce que nos quartiles ne dépendent que de l'ordre des points, pas des valeurs. Une valeur de 0.72 signifie : les deux méthodes s'accordent sur le « gros » (ce qui est très dense ou très vide), mais divergent beaucoup point par point. Concrètement, 63\% des points du Q1 d'ASTRA sont dans le Q1 du density split, et moins de 2\% passent d'un extrême à l'autre.

\paragraph{La figure 5 : pourquoi le density split bouge beaucoup plus qu'ASTRA ?} $\delta_R$ est un contraste de densité « absolu » : si l'Univers est plus grumeleux ($\sigma_8$ plus grand), les régions denses sont plus denses et les vides plus vides, et $\delta_R$ l'enregistre directement. $r$ est un rapport entre voisins galaxies et voisins aléatoires : il est borné entre $-1$ et 1, il sature dans les régions extrêmes, et surtout le voisinage s'adapte à la densité locale. Quand le contraste augmente partout, l'ordre des points change, mais la valeur de $r$ change peu.

\paragraph{Qu'en déduire ? Le clustering d'ASTRA ne se fait-il pas sur $r$ ?} Non : $r$ ne sert qu'à \emph{classer} les points en quartiles. Le clustering est ensuite mesuré sur les \emph{positions} des points de chaque quartile. Et comme les deux méthodes utilisent des quartiles (donc des rangs), aucune des deux n'utilise directement les valeurs $r$ ou $\delta_R$. Ce que la figure 5 nous apprend : (1) la distribution des valeurs de $r$ serait une mauvaise sonde cosmologique (au contraire de la distribution des densités pour le density split) ; (2) l'information d'ASTRA est dans \emph{l'emplacement} de ses environnements et la façon dont ils se regroupent, qui eux réagissent bien à la cosmologie.

\subsection{Section 4 : les corrélations}

\paragraph{RR et les « $5\times$ le nombre de galaxies ».} Pour savoir si des points sont regroupés, on compare le nombre de paires observées au nombre attendu « sans regroupement » dans la même géométrie. On tire pour cela un catalogue de points aléatoires uniformes, cinq fois plus nombreux que les galaxies pour que son propre bruit soit négligeable. $RR$ est le nombre de paires de ces points aléatoires à chaque distance ; $DD$ celui des données, $DR$ les paires mixtes. Comme tous les sous-volumes ont la même forme, $RR$ est le même pour tous et n'est calculé qu'une fois.

\paragraph{« La moyenne des quatre corrélations croisées des quartiles de galaxies est exactement la 2PCF ».} C'est une conséquence du comptage. Chaque paire de galaxies comptée dans la 2PCF a sa première galaxie dans un et un seul des quatre quartiles : elle est comptée une fois dans une des quatre corrélations croisées. Les quatre quartiles ont exactement la même taille et utilisent le même catalogue aléatoire, donc, une fois les comptages normalisés, la moyenne des quatre donne les comptages de la 2PCF ; l'estimateur de Landy--Szalay étant linéaire en ces comptages, la moyenne des quatre corrélations croisées \emph{est} la 2PCF, à un détail de normalisation près (le nombre de paires possibles n'est pas tout à fait le même pour une auto-corrélation et une corrélation croisée), totalement négligeable. (Pour les quartiles de points aléatoires, ce n'est pas le cas : leur moyenne donne la corrélation points aléatoires--galaxies, qui est nulle.)

\paragraph{Pourquoi les auto-corrélations des quartiles extrêmes sont amplifiées.} Les régions très denses et très vides ne sont pas placées au hasard : elles se trouvent là où les grandes fluctuations de densité sont fortes. Elles « tracent » ces grandes fluctuations avec un fort biais (positif pour les amas, négatif pour les vides), et un traceur biaisé a une auto-corrélation amplifiée (en $b^2$). C'est pour cela que ces statistiques sont très sensibles à l'amplitude et à la forme des fluctuations, donc à la cosmologie.

\paragraph{Pourquoi les quadrupôles de Q1 et Q4 ont des signes opposés.} Autour d'un vide, les galaxies s'éloignent (le vide se vide) : dans l'espace des redshifts, le vide paraît étiré le long de la ligne de visée, d'où un quadrupôle positif. Autour d'un amas, les galaxies tombent vers lui : il paraît écrasé, quadrupôle négatif. Ces quadrupôles mesurent la vitesse de croissance des structures.

\paragraph{Pourquoi les signaux ASTRA sont 1.6 fois plus faibles.} ASTRA estime l'environnement d'un point avec seulement $\sim15$ voisins : c'est une estimation bruitée. Un point situé dans une région moyenne peut, par hasard, avoir surtout des voisins galaxies et se retrouver classé en Q4. Chaque quartile ASTRA contient donc une partie de points « mal classés », ce qui dilue le contraste entre quartiles. Le density split moyenne sur une sphère de $10\Mpch$ et est moins bruité. À retenir : un signal plus faible n'est pas forcément moins d'information ; ce qui compte est la réponse aux paramètres comparée au bruit.

\paragraph{Espace réel et espace des redshifts : pourquoi les amplitudes sont plus faibles en espace réel.} Dans l'espace des redshifts, les galaxies qui tombent vers les amas paraissent rapprochées les unes des autres le long de la ligne de visée (effet Kaiser) : les structures semblent plus contrastées, donc la corrélation est plus forte. En espace réel (positions vraies), cet effet n'existe pas : les amplitudes sont plus faibles, et il n'y a pas de quadrupôle.

\paragraph{« Les environnements ASTRA sont moins contrastés mais leur anisotropie est relativement plus forte ».} En valeur absolue, les quadrupôles d'ASTRA sont plus petits (comme ses monopôles). Mais le rapport quadrupôle/monopôle est plus grand pour ASTRA. Une explication plausible : ASTRA définit l'environnement avec les voisins les plus proches, \emph{dans l'espace des redshifts}, à une échelle ($\sim13\Mpch$) où les distorsions sont fortes ; sa sélection d'environnements est donc elle-même déformée par les vitesses, alors que le lissage sphérique du density split moyenne en partie ces déformations. ASTRA est donc relativement plus sensible au champ de vitesses.

\subsection{Section 5 : covariance}

\paragraph{Lire une matrice de corrélation (figure 11).} La matrice est un tableau carré : chaque ligne et chaque colonne correspond à un élément du vecteur de données (un bin de distance, d'un multipôle, d'une statistique). La case $(i,j)$ dit si, d'une réalisation à l'autre, les éléments $i$ et $j$ fluctuent ensemble : rouge (+1) ils montent et descendent ensemble, bleu ($-1$) quand l'un monte l'autre descend, blanc (0) ils sont indépendants. La diagonale vaut toujours 1. Les traits noirs délimitent les blocs : 2PCF, ASTRA Q1, ASTRA Q4, DS Q1, DS Q4. On y lit : (1) les bins voisins d'une même courbe sont très corrélés (bandes rouges le long des diagonales) ; (2) Q4 fluctue avec la 2PCF (régions denses et galaxies vont ensemble) ; (3) Q1 et Q4 sont anti-corrélés (si une réalisation a des amas plus marqués, elle a aussi des vides plus marqués) ; (4) le bloc ASTRA Q1 et le bloc DS Q1 sont presque identiques : leur bruit est en grande partie commun.

\paragraph{Pourquoi un bruit commun rend ASTRA et le density split complémentaires.} Imaginons deux mesures $A = a\,\theta + n$ et $B = b\,\theta + n$, qui contiennent \emph{le même} bruit $n$ mais réagissent différemment au paramètre $\theta$ ($a\ne b$). Séparément, chacune est limitée par $n$. Mais $A - B = (a-b)\,\theta$ ne contient plus de bruit du tout ! En réalité le bruit n'est que partiellement commun, mais l'idée reste : en combinant deux mesures dont le bruit est partagé et la réponse différente, on annule une partie du bruit. C'est pourquoi ASTRA + density split gagne beaucoup plus que deux mesures indépendantes (jusqu'à un facteur 2.4 sur $\omega_b$).

\paragraph{Le facteur de Hartlap.} La covariance est estimée à partir d'un nombre fini de simulations ; elle est donc bruitée. Or on a besoin de son \emph{inverse}, et l'inverse d'une matrice bruitée est en moyenne trop grand : les erreurs prévues seraient trop optimistes. Le facteur de Hartlap corrige exactement ce biais (pour des données gaussiennes). Il dépend du rapport entre la longueur du vecteur de données $N_d$ et le nombre de simulations $N_s$ : proche de 1 si $N_s \gg N_d$, mais 0.31 pour 130 éléments et 192 simulations. D'où l'intérêt de comprimer les vecteurs et d'avoir 1786 boîtes.

\subsection{Section 5 : dérivées}

\paragraph{Le passage sur les paires (vulgarisé).} Pour estimer l'effet d'un paramètre, on compare deux simulations qui ne diffèrent que par ce paramètre et partent des mêmes conditions initiales : elles contiennent les mêmes structures aux mêmes endroits. On les compare région par région (sous-volume par sous-volume) : l'essentiel du « hasard cosmique » (quelles structures se trouvent où) est le même des deux côtés et disparaît dans la différence. Mais il reste un hasard qui n'est pas partagé : le tirage des galaxies dans les halos. Chaque catalogue de galaxies est un tirage indépendant (un halo donné a une galaxie dans un catalogue et pas dans l'autre), et ce bruit-là ne disparaît pas dans la différence.

\paragraph{Terme HOD et terme cosmologique : de quoi parle-t-on ?} La différence entre deux catalogues contient deux effets. Le \emph{terme cosmologique} : l'effet du changement de paramètre cosmologique, c'est la physique qu'on veut mesurer. Le \emph{terme HOD} : l'effet du changement de recette de galaxies entre les deux catalogues ; ce n'est pas de la physique de l'Univers, c'est une différence de modèle de formation des galaxies, qu'on n'a pas choisie. On a estimé le terme HOD grâce à 72 catalogues de la cosmologie de référence avec des HOD différentes (ils disent de combien chaque paramètre HOD change les statistiques), et on l'a comparé au terme cosmologique, les deux en unités de l'erreur statistique : le terme HOD est 10 à 15 fois plus grand.

\paragraph{« Même 8\% laissent un résidu comparable au signal ».} Si le terme HOD vaut 10 fois le signal, et qu'on le corrige avec un modèle qui se trompe de 8\%, l'erreur restante vaut $0.08\times10 = 0.8$ fois le signal : presque autant que ce qu'on cherche à mesurer. La correction ne suffit donc pas.

\paragraph{À quoi servent encore les différences finies ?} (1) Elles ont révélé et mesuré le problème : sans elles on ne saurait pas que l'effet HOD domine d'un facteur 10. (2) Une fois corrigées (même imparfaitement), elles donnent des dérivées obtenues sans réseau de neurones, qui servent à vérifier celles de l'émulateur. (3) Avec plusieurs HOD par cosmologie, elles deviennent la régression de la section 5.4, une deuxième méthode indépendante.

\paragraph{Pourquoi un émulateur (le paragraphe SUNBIRD).} L'émulateur apprend la relation « paramètres cosmologiques + HOD $\to$ statistiques » sur tous les catalogues. Une fois appris, on peut lui demander de changer un paramètre cosmologique \emph{sans toucher à la HOD} : sa dérivée ne contient, par construction, aucun terme HOD. La « différentiation automatique » est simplement la façon dont les bibliothèques de réseaux de neurones calculent exactement les dérivées d'un réseau.

\paragraph{Les choix d'architecture.} Un réseau « entièrement connecté » à 3 couches de 256 neurones (environ 140\,000 poids) : c'est le type de réseau utilisé par SUNBIRD pour le density split, adapté à des fonctions lisses de 20 paramètres et à environ 1000 exemples. L'activation SiLU est lisse ; la ReLU, plus courante, est faite de morceaux de droites, et ses dérivées sont en escalier, ce qui rendrait les dérivées de Fisher bruitées. Adam est l'algorithme d'optimisation standard ; « l'arrêt précoce » arrête l'entraînement quand l'erreur sur des cosmologies mises de côté cesse de baisser, pour éviter le sur-apprentissage. La perte en $\chi^2$ pénalise chaque erreur en proportion de sa significativité statistique. Nous n'avons pas testé d'autres tailles de réseau ; la courbe d'apprentissage indique que la limite vient du bruit des catalogues, pas du réseau.

\paragraph{La précision de l'émulateur est-elle bonne ?} L'erreur médiane est de 0.4 à 0.9 fois l'erreur statistique d'une boîte complète de $(2\,h^{-1}{\rm Gpc})^3$. Cela dépend du volume qu'on veut analyser : pour nos prévisions, faites sur des volumes de $(500\Mpch)^3$ (64 fois plus petits, donc environ 8 fois plus bruités), l'erreur de l'émulateur est environ 10 fois plus petite que l'erreur statistique : négligeable, l'émulateur est parfaitement adapté. Pour un volume de $(2\,h^{-1}{\rm Gpc})^3$ ou un relevé de type DESI, elle devient comparable à l'erreur statistique : une analyse de ces volumes serait limitée par l'émulateur. Il est donc bien adapté aux comparaisons de méthodes et aux petits volumes, pas encore à une analyse finale de DESI.

\paragraph{Le test hors échantillon et la densification.} Prédire correctement des statistiques ne garantit pas de prédire correctement leurs \emph{variations}. Pour le tester, on a réentraîné l'émulateur sans les 14 cosmologies de la grille de dérivées, puis on lui a demandé de prédire les différences entre ces cosmologies. Il retrouve l'effet de $\sigma_8$ et $n_s$ à 10--20\% près, celui de $\omega_b$ à 30--45\%, mais pas celui de $\omega_{\rm cdm}$ aux petites échelles : à $\theta_*$ et $\sigma_8$ fixés, cet effet est plus petit que l'erreur de l'émulateur. On a donc ajouté 240 catalogues proches du point de référence (« densification »), pour que l'émulateur voie beaucoup d'exemples de l'effet de chaque paramètre à HOD similaire. Ça n'a pas suffi : les dérivées variaient encore beaucoup d'un réseau à l'autre (70--170\%). Les dérivées de l'émulateur concordent cependant avec les différences de paires corrigées par le processus gaussien, à 1--13\% selon le paramètre.

\paragraph{L'entraînement sur les différences.} Pour aider l'émulateur à apprendre les \emph{variations}, on ajoute à sa fonction de perte un terme qui le pénalise s'il se trompe sur les différences entre catalogues proches (même cosmologie à HOD différentes, ou cosmologies $\pm$). On a testé en plus un poids qui favorise les petites échelles. Pour juger, on a mis de côté 56 catalogues et vérifié les différences prédites : le terme de différences divise l'erreur par deux sur $\omega_b$ et $\sigma_8$, et l'erreur restante atteint le niveau de bruit des différences mesurées, on ne peut donc pas faire de test plus sévère. Le poids aux petites échelles gagne un peu mais rend certaines prévisions instables : on ne le garde pas.

\paragraph{La figure des dérivées (figure 16).} Elle montre de combien chaque statistique change quand on change un peu chaque paramètre, en unités de l'erreur : une courbe qui monte à 0.5 veut dire que le changement vaut une demi-barre d'erreur. Trois choses à retenir : (1) $\omega_{\rm cdm}$ et $n_s$ changent la \emph{forme} de la corrélation (surtout aux petites échelles), de la même façon pour ASTRA et le density split ; (2) les vides (Q1) réagissent dans le sens opposé aux amas (Q4) et à la 2PCF. C'est crucial : un changement de biais des galaxies (effet HOD) change l'amplitude de toutes les courbes dans le même sens, alors qu'un paramètre cosmologique fait monter les unes et descendre les autres ; en les combinant, on distingue les deux, on « lève la dégénérescence » ; (3) à HOD fixée, $\sigma_8$ agit surtout sur les quadrupôles, c'est-à-dire via les vitesses.

\paragraph{Les figures 17 et 18 et le panneau vide.} C'était un bug : la figure traçait la dérivée multipliée par la valeur du paramètre. Or $w_a=0$ au point de référence, d'où une courbe nulle (et un signe inversé pour $w_0=-1$). La figure trace maintenant la réponse à un petit pas fixé de chaque paramètre, comme la figure 16.

\paragraph{La régression « sous-volume par sous-volume » et le bootstrap.} On ajuste la régression séparément dans chaque sous-volume : comme les deux cosmologies partent des mêmes conditions initiales, un même sous-volume contient les mêmes structures dans les deux, et leur différence est débarrassée du hasard cosmique. Le bootstrap estime l'incertitude : on refait l'ajustement sur des tirages aléatoires (avec remise) des sous-volumes, et la dispersion des résultats donne l'erreur.

\paragraph{L'espace HOD réduit, en termes statistiques généraux.} On ajuste un modèle à 12 paramètres avec environ 18 points : presque autant d'inconnues que de données. Le modèle peut alors passer par tous les points, mais il ne prédit rien de nouveau : c'est le sur-apprentissage. On le détecte par validation croisée « leave-one-out » : on retire un point, on ajuste sur les autres, on prédit le point retiré, et on recommence pour chaque point. Ici l'erreur de prédiction valait 0.8--1.0 fois la dispersion des données, c'est-à-dire guère mieux que de prédire la moyenne. Le remède classique est de réduire le nombre de directions : une décomposition en valeurs singulières (proche d'une analyse en composantes principales) de l'effet de la HOD montre qu'une seule combinaison des 12 paramètres compte vraiment (valeurs singulières $1, 0.05, 0.01, \dots$). En ne gardant que 3 directions, l'ajustement devient prédictif.

\paragraph{Pourquoi on ajoute l'erreur de l'émulateur à la covariance.} Les mesures ont leur propre bruit (covariance mesurée), et le modèle (l'émulateur) a sa propre erreur. Les deux sont indépendants ; or des incertitudes indépendantes s'additionnent en variance. Ignorer l'erreur de l'émulateur rendrait les prévisions trop optimistes.

\paragraph{Fisher et compression en deux mots.} $D$ dit de combien les données bougent quand chaque paramètre bouge ; $C^{-1}$ pondère ces mouvements par l'inverse du bruit ; $P_h$ ajoute ce qu'on sait déjà des paramètres HOD. La matrice de Fisher $F = D^TC^{-1}D + P_h$ donne, une fois inversée, les incertitudes attendues. La compression $\boldsymbol t = W^T\boldsymbol d$ remplace le long vecteur de données par une combinaison par paramètre, la plus sensible possible à ce paramètre. Si la compression est imparfaite, on perd de l'information, mais on ne peut pas \emph{fausser} la prévision dans le sens optimiste : on calcule l'information de $\boldsymbol t$ avec sa vraie covariance, et une version comprimée des données ne peut jamais contenir plus d'information que les données complètes. Au pire, les erreurs prévues sont un peu plus grandes.

\paragraph{« La marginalisation sur la HOD dégrade les contraintes de moins de 30\% ».} Certains paramètres HOD imitent l'effet de paramètres cosmologiques (par exemple, le biais des galaxies change l'amplitude comme $\sigma_8$). Si on laisse la HOD libre (« marginaliser »), une partie du signal cosmologique peut être expliquée par la HOD, et les incertitudes grandissent. Une dégradation de moins de 30\% signifie que nos statistiques savent bien distinguer la cosmologie de la formation des galaxies ; $n_s$ est le plus touché parce que son effet ressemble le plus à un effet HOD.

\subsection{Section 6 : le tableau de l'évolution de l'analyse (tableau 5)}

Ce tableau résume l'histoire du projet en cinq corrections, chacune motivée par l'échec de la précédente :
\begin{enumerate}
\item \textbf{Moyenner des fichiers HOD} $\to$ \textbf{paires de HOD proches} : les 500 fichiers d'une cosmologie sont 500 modèles de galaxies différents (55\% d'écart), pas des répétitions ; on ne peut pas les moyenner pour mesurer un effet de 1\%.
\item \textbf{Différences finies} $\to$ \textbf{émulateur} : même en choisissant des HOD proches, l'effet HOD reste 10 fois plus grand que l'effet cosmologique ; l'émulateur dérive à HOD exactement fixée.
\item \textbf{Émulateur} $\to$ \textbf{émulateur densifié et entraîné sur les différences} : le premier émulateur n'apprenait pas bien les petites variations (dérivées instables) ; on lui a donné plus d'exemples locaux et une perte qui porte sur les variations.
\item \textbf{192 sous-volumes d'une boîte} $\to$ \textbf{1786 boîtes indépendantes} : la covariance tirée d'une seule simulation était biaisée ; avec des boîtes indépendantes, la petite avance d'ASTRA sur $\sigma_8$ disparaît.
\item \textbf{Quijote à 6 paramètres} $\to$ \textbf{sous-ensembles fiables} : avec des traceurs rares, le bruit des dérivées domine pour les paramètres dégénérés ; on ne garde que les combinaisons où deux estimateurs de Fisher indépendants sont d'accord.
\end{enumerate}
La colonne de droite montre l'effet de chaque correction sur la comparaison ASTRA / density split : c'est ce qui permet de dire quelles conclusions sont robustes (information comparable, complémentarité) et lesquelles dépendaient des choix d'analyse (le classement fin paramètre par paramètre).
